\chapter{Trabajos relacionados}
\label{chap:relacionados}

\section{Alcance de la revisión}

La revisión se centra en trabajos que ayudan a responder la pregunta experimental del TFG. Se consideran cinco líneas. La primera estudia el diagnóstico automático a partir de señales de ECG. La segunda utiliza el trazado como imagen. La tercera genera datos sintéticos. La cuarta desarrolla modelos multimodales para imágenes electrocardiográficas. La quinta analiza aprendizaje en contexto con demostraciones visuales. Las dos últimas líneas reciben mayor atención porque constituyen la dirección experimental principal del proyecto.

La búsqueda bibliográfica se actualizó hasta el 22 de junio de 2026. Se priorizaron artículos revisados por pares, conjuntos publicados con documentación técnica, guías clínicas y documentación oficial de modelos. Los resultados no se comparan como si procedieran del mismo banco de pruebas. Las bases de datos, las etiquetas y las métricas difieren entre estudios.

Esta precaución es necesaria porque una exactitud alta en una tarea binaria no puede compararse directamente con una clasificación multietiqueta. Tampoco es equivalente evaluar imágenes renderizadas desde señales limpias que fotografías de documentos clínicos. La revisión se utiliza para identificar decisiones de diseño y huecos metodológicos.

\section{Aprendizaje profundo sobre señales electrocardiográficas}

El análisis automático moderno se desarrolló principalmente sobre señales temporales. Las convoluciones unidimensionales permiten aprender filtros directamente sobre las muestras y evitan diseñar manualmente todas las características. Las redes profundas amplían el campo receptivo hasta integrar información distribuida a lo largo de varios segundos.

Ribeiro y colaboradores entrenaron una red profunda con más de dos millones de ECG de doce derivaciones procedentes de una red brasileña de telemedicina \cite{ribeiro2020}. El sistema reconocía seis alteraciones y fue comparado con médicos en un subconjunto anotado. El trabajo mostró que una gran colección clínica puede sostener un clasificador temporal de alto rendimiento.

Su escala también señala una limitación para enfermedades poco frecuentes. Reunir millones de registros con etiquetas fiables no es una estrategia disponible para todos los problemas. Además, las etiquetas de rutina pueden reflejar criterios locales, versiones de software y variabilidad entre observadores.

PTB XL proporcionó un punto de referencia abierto con 21 837 ECG de doce derivaciones y 71 enunciados potencialmente multietiqueta \cite{wagner2020}. Sus particiones recomendadas facilitan comparaciones reproducibles. MIMIC IV ECG amplió la escala hasta aproximadamente 800 000 registros vinculables con información clínica \cite{gow2023}. Estos recursos resultan valiosos para preentrenamiento y validación, aunque no garantizan una cantidad suficiente de cada condición rara.

La literatura sobre señal establece dos referencias para el presente trabajo. Primero, una arquitectura supervisada puede aprender morfología electrocardiográfica cuando dispone de datos adecuados. Segundo, el rendimiento depende de la correspondencia entre el dominio de entrenamiento y la población objetivo. El simulador se utiliza para estudiar el régimen anterior a una gran colección clínica, no para competir con sistemas entrenados sobre millones de registros.

\section{Análisis directo de imágenes de ECG}

El análisis de imagen aborda situaciones en las que el trazado está disponible como documento. Gliner y colaboradores compararon redes que recibían señales digitales y redes que recibían imágenes de doce derivaciones \cite{gliner2020}. Su estudio incluyó ocho condiciones para el modelo temporal y una evaluación visual centrada en fibrilación auricular. También aplicó transformaciones que imitaban capturas con teléfono.

Sangha y colaboradores desarrollaron una CNN multietiqueta a partir de 2 228 236 registros brasileños convertidos en imagen \cite{sangha2022}. El modelo predijo seis etiquetas clínicas y se evaluó en una cohorte externa alemana y en ECG impresos. El estudio demostró que la representación visual puede generalizar cuando el entrenamiento incorpora escala, diversidad y variación de formato.

La principal diferencia con este TFG es la disponibilidad de datos. Sangha y colaboradores estudian aprendizaje supervisado a gran escala. El presente trabajo pregunta qué ocurre antes de alcanzar esa escala y compara el entrenamiento de una CNN con el uso inmediato de demostraciones en un modelo abierto.

La imagen puede analizarse directamente o convertirse de nuevo a señal. Wu y colaboradores desarrollaron un sistema automático de digitalización y evaluaron la correlación entre la señal reconstruida y la original \cite{wu2022digitization}. Los resultados fueron altos en formatos favorables y disminuyeron cuando había solapamiento entre derivaciones. Esta variabilidad muestra que la digitalización añade sus propios fallos.

El trabajo de Gliner y colaboradores de 2023 se centra en la transferencia desde imágenes limpias a fotografías y formatos no observados \cite{gliner2023}. Su arquitectura adversarial utiliza muestras objetivo sin etiqueta para reducir la dependencia del estilo de adquisición. Esta línea es directamente relevante para la futura transición del simulador a datos reales.

\section{Datos sintéticos para ECG}

La generación sintética aparece en varios niveles de fidelidad. Los modelos gaussianos representan ondas mediante funciones paramétricas. Su principal ventaja es la interpretabilidad. Cada cambio puede relacionarse con amplitud, posición o anchura. Sayadi y colaboradores emplearon un modelo de este tipo dentro de un sistema dinámico para generación y filtrado \cite{sayadi2010}.

Los modelos electrofisiológicos simulan la propagación eléctrica sobre geometrías cardíacas. MedalCare XL contiene 16 900 ECG sintéticos de doce derivaciones para un grupo sano y siete patologías \cite{gillette2023}. La colección aporta señales con una base mecanística y permite estudiar condiciones difíciles de equilibrar en datos clínicos.

También existen modelos generativos aprendidos. Los modelos de difusión condicionados pueden producir señales asociadas a enunciados diagnósticos \cite{alcaraz2023}. Su flexibilidad es mayor que la de un simulador manual, pero la relación entre cada parámetro y la morfología final resulta menos directa. Además, la calidad debe evaluarse mediante diversidad, fidelidad clínica y utilidad para tareas posteriores.

El simulador del proyecto se mantiene en una versión simple. No intenta reproducir un torso, un corazón tridimensional o una distribución clínica completa. Genera un latido para estudiar tres hallazgos localizados. Esta decisión reduce el realismo y mejora la capacidad de auditar las reglas.

La literatura sintética también advierte de un riesgo. Un conjunto generado puede amplificar los supuestos del modelo de simulación. Si la variabilidad no incluye ruido, escalas y formas intermedias, el clasificador aprenderá una frontera más limpia que la que existe en práctica clínica. La evaluación real sigue siendo necesaria.

\section{Conjuntos que integran señal, imagen y texto}

Los modelos multimodales necesitan datos alineados. Una imagen sin informe permite clasificación supervisada, pero no ofrece una relación rica entre morfología y lenguaje. Un informe sin trazado tampoco permite comprobar si la explicación se apoya en evidencia visual.

MEETI se publicó en 2026 como una extensión de MIMIC IV ECG \cite{zhang2026meeti}. El recurso sincroniza señales, imágenes renderizadas, parámetros por latido e interpretaciones textuales. Cubre hasta 784 680 registros de 160 597 pacientes. Las imágenes se generan con una escala estándar y las interpretaciones detalladas se producen mediante un modelo de lenguaje a partir de informes y parámetros.

MEETI reduce la barrera de entrada para estudiar alineamiento multimodal a gran escala. También requiere una lectura crítica. Parte del texto es generado por un modelo y no constituye una nueva anotación clínica independiente. El recurso resulta adecuado para preentrenamiento, recuperación y generación de informes, pero la evaluación diagnóstica debe conservar etiquetas y conjuntos de referencia validados.

El proyecto no utiliza por ahora MEETI porque su tarea se limita a un latido sintético y tres hallazgos. Sin embargo, este conjunto ofrece una ruta futura para ampliar la representación hacia doce derivaciones y para comparar imagen, señal y texto sobre los mismos registros.

\section{Modelos multimodales especializados en ECG}

Los modelos generalistas de visión y lenguaje han sido entrenados principalmente con imágenes naturales, documentos y contenido web. Un trazado electrocardiográfico posee una estructura diferente. La información se codifica en deflexiones delgadas, relaciones temporales y comparaciones entre derivaciones. Esta distancia de dominio puede limitar la interpretación directa y motiva dos aproximaciones complementarias: adaptar un modelo de gran tamaño al dominio médico, o aprovechar las capacidades de razonamiento visual de un modelo generalista sin modificar sus pesos. Ambas vías se exploran en este TFG.

Liu y colaboradores introdujeron ECGInstruct, PULSE y ECGBench \cite{liu2026}. ECGInstruct contiene más de un millón de pares de imagen y texto que cubren reconocimiento de características, ritmo, morfología y generación de informes. PULSE es un modelo de siete mil millones de parámetros adaptado a ECG. ECGBench evalúa cuatro tareas en nueve conjuntos e incluye imágenes sintéticas y reales.

Los autores informan que PULSE supera a modelos multimodales generalistas en su banco de pruebas \cite{liu2026}. Este resultado apoya la idea de que el conocimiento visual específico importa. Al mismo tiempo, PULSE se ha entrenado con una colección masiva de instrucciones. Su existencia confirma que la adaptación al dominio mejora el rendimiento, pero no responde a la misma pregunta que este TFG. El presente proyecto investiga qué rendimiento se puede obtener \emph{sin ajuste} de pesos, utilizando un modelo generalista abierto cuando solo existen unos pocos ejemplos de una condición rara como Brugada.

PULSE constituye una referencia futura importante. Si se incorpora al protocolo, permitiría separar tres niveles de especialización. El primero sería una CNN entrenada de forma supervisada con pocas muestras. El segundo sería un modelo generalista mediante aprendizaje en contexto. El tercero sería un modelo multimodal ya adaptado a ECG. La comparación entre estos tres niveles mediría el valor incremental de cada etapa de especialización para una tarea de baja prevalencia.

La aparición de PULSE modifica el estado del arte, pero no elimina la pregunta de investigación. Un modelo especializado puede no estar disponible para una modalidad rara, una definición local o un nuevo protocolo de anotación. El aprendizaje en contexto sigue siendo relevante durante las fases tempranas de formulación y recogida de datos, especialmente cuando el coste de entrenamiento especializado es prohibitivo y se necesita una evaluación rápida de viabilidad.

\section{Aprendizaje multimodal en contexto}

El aprendizaje en contexto (\emph{in-context learning}, ICL) modifica el ciclo habitual de entrenamiento supervisado. En lugar de ajustar los pesos de una red mediante descenso por gradiente, los ejemplos se colocan en la ventana de contexto del modelo durante la inferencia. Cada demostración consiste en una imagen y su etiqueta o descripción textual. El modelo produce una predicción para la imagen de consulta sin que sus pesos hayan cambiado. Esta propiedad es especialmente valiosa cuando los datos disponibles se reducen a unas pocas decenas de muestras, situación habitual en condiciones cardíacas de baja prevalencia como el síndrome de Brugada.

Flamingo estableció la arquitectura de referencia para esta capacidad. El modelo recibe secuencias intercaladas de imágenes y texto y resuelve tareas con demostraciones de ejemplo \cite{alayrac2022}. Flamingo demostró que un modelo preentrenado de gran escala puede adaptarse a tareas nuevas durante la inferencia cuando se le proporcionan ejemplos visuales y textuales relevantes. Este resultado impulsó el uso de modelos multimodales como sistemas flexibles que no requieren un ciclo de ajuste dedicado para cada tarea.

Sin embargo, los estudios posteriores demuestran que esta capacidad es frágil y que la contribución de cada componente de las demostraciones no es homogénea. Chen y colaboradores evaluaron si los modelos utilizaban realmente la parte visual de las demostraciones \cite{chen2024}. Encontraron una dependencia dominante del contenido textual en varios bancos de pruebas estándar. La selección de ejemplos basada en similitud visual seguía aportando valor marginal, pero la imagen dentro de la demostración podía tener menos influencia de la esperada. Este hallazgo plantea una pregunta directa para el presente trabajo: cuando las demostraciones contienen imágenes de ECG, ¿el modelo las procesa visualmente o se apoya solo en las descripciones textuales?

Qin y colaboradores estudiaron la selección, el orden y la construcción de los ejemplos \cite{zhang2024mmicl}. Sus resultados muestran que la recuperación multimodal, la posición de la instrucción dentro del contexto y la diversidad de las demostraciones afectan al rendimiento de forma significativa. Esto impide tratar \(K\) como la única variable relevante. Dos contextos con el mismo número de ejemplos pueden contener información muy diferente y producir resultados dispares.

Estas observaciones han influido directamente en el diseño experimental del repositorio. Cada valor positivo de \(K\) se repite con tres semillas para cuantificar la variabilidad asociada a la selección de ejemplos. Las variantes con respuestas permutadas sirven como control para detectar si el modelo está memorizando la correspondencia etiqueta-posición en lugar de razonar visualmente. Las variantes sin imagen miden la contribución del canal visual. Las instrucciones breves y las descripciones morfológicas permiten comparar el efecto del conocimiento textual explícito frente al razonamiento visual implícito. En conjunto, estos controles permiten atribuir el rendimiento a cada componente del contexto y evitar conclusiones espurias sobre la capacidad visual del modelo.

Además, la aplicación de ICL a dominios médicos especializados plantea retos adicionales que la literatura general no ha abordado en profundidad. La morfología de un ECG carece de los objetos semánticos familiares que abundan en los bancos naturales. El modelo debe interpretar deflexiones de pocos milímetros, elevaciones del segmento ST y relaciones temporales entre ondas. No es evidente que la representación visual preentrenada sobre fotografías y documentos sea suficiente para capturar estas diferencias sutiles. El presente trabajo es, hasta donde sabemos, uno de los primeros en aplicar ICL multimodal a imágenes de ECG de un solo latido con condiciones raras.

\section{Gemma 4 y MedGemma}

Gemma 4 es una familia abierta y multimodal publicada por Google DeepMind en 2026. Sus variantes admiten texto e imagen, contexto extenso y resolución visual variable \cite{gemma4card}. El presupuesto de tokens visuales permite estudiar un compromiso entre detalle y coste. La posibilidad de ejecución local facilita experimentos reproducibles y reduce la necesidad de enviar imágenes a un servicio externo.

El proyecto selecciona Gemma 4 como modelo generalista principal para la línea VLM/ICL. Esta elección responde a criterios prácticos y metodológicos. En primer lugar, Gemma 4 ofrece pesos abiertos, lo cual permite una reproducibilidad completa y evita la dependencia de un servicio propietario cuyo comportamiento puede cambiar entre versiones. En segundo lugar, su ventana de contexto admite varias imágenes simultáneas, requisito imprescindible para implementar demostraciones de tipo \emph{few-shot} con \(K\) ejemplos visuales. En tercer lugar, la compatibilidad con servidores de inferencia locales permite ejecutar cientos de evaluaciones sin coste por petición.

La elección de Gemma 4 no se justifica mediante una afirmación previa de superioridad en ECG. No existen, hasta donde sabemos, estudios publicados que evalúen este modelo sobre trazados electrocardiográficos. El rendimiento deberá establecerse mediante los experimentos del capítulo~\ref{chap:vlm}. Precisamente esta ausencia de evaluación previa refuerza la motivación del TFG: determinar si un modelo generalista abierto puede interpretar imágenes de ECG sin ajuste cuando se le proporcionan unas pocas demostraciones.

MedGemma 1.5 se incluye como referencia médica complementaria. El modelo ha recibido adaptación con diferentes modalidades clínicas, incluyendo radiología y dermatología \cite{medgemma15}. La documentación disponible no sitúa el ECG entre sus modalidades principales, por lo que su uso en este dominio requiere validación específica. La comparación entre Gemma 4 y MedGemma permitirá responder a una pregunta relevante: ¿el preentrenamiento médico general aporta una ventaja sobre un modelo generalista más reciente cuando la tarea consiste en interpretar un trazado de ECG de un solo latido? Si la ventaja es nula o marginal, se concluirá que la adaptación médica general no transfiere automáticamente al dominio electrocardiográfico, lo cual reforzaría la relevancia de modelos especializados como PULSE.

Las versiones exactas deberán registrarse porque Gemma 4 y MedGemma pertenecen a generaciones base distintas. También deben conservarse la resolución, el presupuesto visual, la plantilla de conversación y la configuración de inferencia. Sin esa información la comparación no sería reproducible. El protocolo del repositorio registra automáticamente estos metadatos en cada ejecución experimental.

\section{Síntesis comparativa}

La tabla \ref{tab:related_work} resume los trabajos más relacionados. Las cifras indican la escala del estudio original y no deben interpretarse como tamaños directamente comparables. Algunas colecciones contienen señales, otras imágenes renderizadas y otras pares multimodales.

\begin{table}[H]
  \centering
  \caption{Síntesis de trabajos relacionados}
  \label{tab:related_work}
  \scriptsize
  \begin{tabularx}{\textwidth}{>{\raggedright\arraybackslash}p{0.14\textwidth} >{\raggedright\arraybackslash}p{0.13\textwidth} >{\raggedright\arraybackslash}p{0.18\textwidth} >{\raggedright\arraybackslash}p{0.19\textwidth} X}
    \toprule
    Trabajo & Entrada & Escala & Aportación & Relación con este TFG \\
    \midrule
    Ribeiro et al. \cite{ribeiro2020} & Señal de doce derivaciones & Más de dos millones de ECG & Clasificación profunda de seis alteraciones & Referencia supervisada a gran escala \\
    Wagner et al. \cite{wagner2020} & Señal de doce derivaciones & 21 837 ECG & Conjunto abierto multietiqueta con particiones recomendadas & Ejemplo de recurso real para transferencia futura \\
    Gliner et al. \cite{gliner2020} & Señal e imagen & 41 830 registros recopilados & Comparación entre entrada digital y visual con artefactos de captura & Evidencia de viabilidad del ECG como imagen \\
    Sangha et al. \cite{sangha2022} & Señal e imagen & 2 228 236 ECG de entrenamiento & Diagnóstico multietiqueta con validación externa e impresa & Referencia visual supervisada a gran escala \\
    Wu et al. \cite{wu2022digitization} & Imagen de papel & 515 registros de validación & Reconstrucción automática de señales desde documentos & Alternativa al análisis directo de imagen \\
    Gliner et al. \cite{gliner2023} & Imagen limpia y fotografía & 79 226 imágenes etiquetadas y 39 613 imágenes objetivo sin etiqueta & Adaptación adversarial entre formatos & Base para la transferencia sintético real \\
    Gillette et al. \cite{gillette2023} & Señal sintética de doce derivaciones & 16 900 ECG & Simulación electrofisiológica de ocho grupos & Contraste con el simulador paramétrico simple \\
    Liu et al. \cite{liu2026} & Imagen y texto & Más de un millón de instrucciones & PULSE, ECGInstruct y ECGBench & Referencia especializada para modelos multimodales \\
    Zhang et al. \cite{zhang2026meeti} & Señal, imagen, parámetros y texto & Hasta 784 680 registros & Recurso multimodal alineado basado en MIMIC IV ECG & Ruta de ampliación hacia datos clínicos multimodales \\
    \bottomrule
  \end{tabularx}
\end{table}

\section{Hueco de investigación}

La literatura revisada demuestra que el ECG puede analizarse como señal y como imagen con resultados clínicamente útiles, siempre que se disponga de una cantidad suficiente de datos etiquetados. Cuando la escala de entrenamiento se reduce a millones o cientos de miles de muestras, los modelos supervisados alcanzan un rendimiento competitivo con expertos. Sin embargo, para condiciones de baja prevalencia como el síndrome de Brugada, que puede representar menos del uno por ciento de un archivo clínico, ninguno de los trabajos revisados aborda la situación en la que el presupuesto total de ejemplos etiquetados se limita a unas pocas decenas.

La línea CNN de este TFG documenta precisamente esta situación. Los resultados del capítulo experimental muestran que una red convolucional entrenada con muy pocas muestras alcanza un rendimiento limitado y pierde capacidad diagnóstica al transferirse a un dominio real. Este resultado no es un defecto del diseño, sino la consecuencia esperable de un entrenamiento supervisado sin datos suficientes. La línea CNN sirve, por tanto, como base para medir el coste de la especialización cuando los datos son escasos.

Frente a esta limitación, los modelos de visión y lenguaje abren una vía alternativa. En lugar de modificar pesos mediante descenso por gradiente, un VLM recibe los mismos \(K\) ejemplos como demostraciones dentro de su ventana de contexto y produce una predicción sin entrenamiento. Esta aproximación posee tres ventajas potenciales en el régimen de pocos datos. Primera, no requiere un ciclo de entrenamiento y, por tanto, elimina el riesgo de sobreajuste a un conjunto diminuto. Segunda, el modelo puede integrar conocimiento morfológico previo adquirido durante su preentrenamiento masivo. Tercera, el coste computacional de la inferencia es predecible y no depende del número de épocas.

Hasta donde sabemos, ningún trabajo publicado ha realizado la comparación directa que propone este TFG: dado un presupuesto fijo de \(K\) imágenes de ECG etiquetadas con una condición rara, ¿es más eficaz utilizar esas \(K\) muestras para entrenar una CNN desde cero o para construir un contexto de demostraciones en un VLM abierto? Los trabajos de Ribeiro et al. \cite{ribeiro2020} y Sangha et al. \cite{sangha2022} estudian el entrenamiento supervisado a gran escala. Los trabajos de Alayrac et al. \cite{alayrac2022}, Chen et al. \cite{chen2024} y Qin et al. \cite{zhang2024mmicl} estudian el aprendizaje en contexto multimodal sobre tareas naturales. Liu et al. \cite{liu2026} especializan un modelo multimodal con más de un millón de instrucciones. Ninguna de estas líneas responde a la comparación planteada con el mismo presupuesto \(K\) sobre imágenes de ECG y, específicamente, sobre una condición rara.

El TFG aborda ese hueco mediante un diseño experimental que impone un presupuesto común de ejemplos a ambos enfoques. La CNN utiliza las \(K\) muestras para modificar sus pesos. Gemma 4 y MedGemma las utilizan dentro del contexto de inferencia. Los controles de imagen eliminada y etiquetas permutadas permiten comprobar si el modelo multimodal aprovecha de verdad la correspondencia visual. El simulador proporciona una tarea conocida y verificable antes de introducir la variabilidad del dominio clínico.

La contribución esperada no es una nueva arquitectura ni un nuevo modelo entrenado. Es una caracterización experimental del modo de uso. La línea CNN, ya cerrada en esta memoria, responde a la primera mitad de la pregunta: qué ocurre cuando se entrena un modelo visual específico con muy pocos ejemplos y se intenta transferirlo a un dominio real. Los resultados muestran que el entrenamiento con datos insuficientes produce modelos frágiles. La línea VLM/ICL deberá completar la segunda mitad: comprobar si el mismo presupuesto \(K\), usado como contexto en lugar de como entrenamiento, mejora la relación entre rendimiento, coste y trazabilidad. Si el VLM supera a la CNN en este escenario, se podrá sostener que, para condiciones raras del ECG, el aprendizaje en contexto con un modelo abierto es una alternativa más eficaz que el entrenamiento supervisado tradicional con datos escasos.
